\documentclass[a4paper]{article}
\usepackage[pdftex]{graphicx}
\usepackage[utf8]{inputenc}
\usepackage{enumerate}
\usepackage{amssymb}
\usepackage{icomma}
\usepackage{tikz}
\usepackage{href-ul}
\hypersetup{
	colorlinks=true,
	linkcolor=blue,
	urlcolor=blue}
\usepackage{geometry}
\geometry{a4paper, top=15mm, left=15mm, right=15mm, bottom=15mm,
	headsep=10mm, footskip=12mm}

\begin{document}
\begin{tikzpicture}(10,3)
	\draw[ultra thick](2,0) --(10,0) -- (10,1.3) --(2,1.3) -- (2,0);
	\draw[fill=black](2,0)-- (0,.6) -- (2,.6) -- (2,0);
	\node at (6,.6) {\large Lösungsblatt von \href{https://www.okuyakl.de}{www.okuyakl.de}};
\end{tikzpicture}
\vspace{0.5 cm}

\noindent{\bf Aufgabe 1.}\\
	\begin{minipage}{0.5\textwidth}
	\renewcommand{\arraystretch}{2}
	\begin{tabular}{|c|c|c|}
		\hline
		& stetig & differenzierbar \\
		\hline
		a) 	&  ja & nein\\
		\hline
		b) 	& ja  & ja\\
		\hline
		c) 	& nein  & nein\\
		\hline
		d) 	& ja  & ja\\
		\hline
	\end{tabular}
\end{minipage}
\hfill
\begin{minipage}{0.5\textwidth}
{\bf e)} Bei $x =-1$ nicht stetig,\\ bei $x=3$ stetig, nicht differenzierbar.\\
\vspace{0.5 cm}

\noindent
{\bf f)} Bei $x=0$ stetig differenzierbar,\\ bei $x=2$ stetig, nicht differenzierbar.
	
\end{minipage}
\vspace{0.5 cm}

\noindent{\bf Aufgabe 2.}\\

\noindent
a) Für $a=3$ stetig, nicht differenzierbar. 
\vspace{0.5 cm}

\noindent
b) Für $t=-2$ stetig und differenzierbar.
\vspace{0.5 cm}

\noindent
c) Für $b={1\over 3}$ stetig, nicht differenzierbar.
\vspace{0.5 cm}

\noindent
d) Für $c=2$ stetig und differenzierbar.
\vspace{0.5 cm}


\noindent{\bf Aufgabe 3.}\\
\vspace{0.5 cm}

\noindent
Diese Betragsfunktionen sind alle stetig.
\vspace{0.5 cm}

\noindent
a) Bei $x=-5$ nicht differenzierbar.
\vspace{0.5 cm}

\noindent
b) Bei $x=t$ nicht differenzierbar.
\vspace{0.5 cm}

\noindent
c) Bei $x=8$ nicht differenzierbar.
\vspace{0.5 cm}

\noindent
d) Bei $x=2$ nicht differenzierbar.
\vspace{0.5 cm}

\noindent
e) Überall differenzierbar.
\vspace{0.5 cm}

\noindent
f) Bei $x=-2$ und $x=3$ nicht differenzierbar.
\vspace{0.5 cm}

\noindent
g) Überall differenzierbar. 
\vspace{0.5 cm}

\noindent
h) Überall differenzierbar. 


\begin{center}
	\includegraphics[width=7 cm]{../../viecher/endcomic.pdf}
	
	Hier geht es zurück zum \href{https://www.okuyakl.de/math/m11dbkL072/aa072.pdf}{Aufgabenblatt}
\end{center}

\end{document}

